\documentclass[11pt]{article}
\usepackage[letterpaper,margin=0.92in]{geometry}
\usepackage{amsmath,amssymb,amsthm}
\usepackage[colorlinks=true,urlcolor=blue,linkcolor=blue,citecolor=blue]{hyperref}
\hypersetup{pdftitle={The Erdos-Dushnik-Miller theorem and countable unions},pdfauthor={Wentao Li}}
\newtheorem{theorem}{Theorem}
\newtheorem{lemma}[theorem]{Lemma}
\newtheorem{corollary}[theorem]{Corollary}
\newcommand{\Ical}{\mathcal I}
\newcommand{\Fix}{\operatorname{Fix}}
\newcommand{\Sym}{\operatorname{Sym}}
\newcommand{\BG}{\mathrm{EDM}_{\mathrm{BG}}}
\newcommand{\ET}{\mathrm{EDM}_{\mathrm T}}
\title{The Erdős--Dushnik--Miller theorem and countable unions}
\author{Wentao Li}
\date{Research draft prepared with assistance from OpenAI GPT-6.2 Astra}
\begin{document}
\maketitle
\begin{abstract}
Two choiceless Erdős--Dushnik--Miller principles imply that one can select from each member of a countably indexed family of countably infinite sets. Howard and Tachtsis ask whether they also make the union of such a family countable. We give a negative answer in ZFA and pure ZF. In a permutation model the atoms are the union of countably many internally countable rows. On each row all permutations moving a density-zero set are allowed. A support argument using the small-index theorem for the infinite symmetric group reduces a difficult graph to a family of copies of individual rows. Ramsey's theorem on their indices then gives the required internal homogeneous set. The graph principle transfers to pure sets by restricting to vertex sets without uncountable well-orderable subsets and treating the remaining case on a regular first uncountable ordinal.
\end{abstract}

\section{The question}
A set is \emph{countable} if it injects into $\omega$. The assertion $\BG$ says that any graph on an uncountable vertex set has a countably infinite independent set or an uncountable clique. The assertion $\ET$ weakens the first alternative to an infinite independent set. The countable union theorem $\mathrm{CUT}$ says that a denumerable union of denumerable sets is denumerable. Howard and Tachtsis \cite[Lemma~5.6]{ht} show that even $\ET$ implies choice from a countable family of denumerable sets. They ask if either principle implies $\mathrm{CUT}$ \cite[Section~8, Question~3]{ht}.
\begin{theorem}\label{main}
Relative to the consistency of ZF, there is a model of pure ZF in which $\aleph_1$ is regular, $\BG$ holds, and a countably indexed disjoint union of denumerable sets is uncountable. The same conjunction has a permutation model with atoms.
\end{theorem}
Since $\BG$ implies $\ET$, this answers both versions of the question. The stronger partition principle requiring a clique equinumerous with the whole vertex set is outside the statement.

\section{The row group and its supports}
Work in a ground universe with choice and atoms $A=\bigsqcup_{n<\omega}A_n$, each row supplied \emph{externally} with an enumeration by $\omega$. Call a subset of a row \emph{small} if its set of positions has asymptotic density zero. We use only that the small sets form a proper ideal, contain the finite sets, and form a tall ideal: each infinite set contains an infinite small subset. For density zero the last fact follows by choosing an increasing subsequence with its $k$th term at least $k^2$.

Let $G_n$ contain all permutations of $A_n$ whose moved points form a small set, and put $G=\prod_nG_n$, the full product. A permutation in $G_n$ preserves the ideal since it changes any set only inside its small set of moved points. A support is an $E\subseteq A$ such that $E_n=E\cap A_n$ is small except at finitely many rows. Such exceptional sections may always be enlarged to whole rows. These supports form an invariant ideal. Use the filter generated by their pointwise stabilizers to form the usual hereditary permutation model. All external choices in the argument below occur in the ground universe. A ground set of internal objects is internal when it has an allowed support, and it is internally well-orderable if a common support fixes all its members individually.

The indexed sequence of rows has empty support. Each row is denumerable internally, as its ground enumeration is supported by the whole row. An internal injection of $A$ into $\omega$ would have a support whose stabilizer fixes every atom. Outside its finitely many exceptional rows there are transpositions of two atoms outside the support. This proves failure of $\mathrm{CUT}$.

Fix for now a support $E$ with exceptional rows filled out and put $K=\Fix_G(E)$. A relative support $S$ for $x$ means that $\Fix_G(E\cup S)$ fixes $x$. We use $S$ outside $E$. In each nonexceptional row the residual row group on $A_n\setminus E_n$ consists of all permutations with small moved set in the restricted ideal, which is tall and proper.

\begin{lemma}\label{intersect}
Suppose $S,T$ are relative supports of $x$. At each residual row suppose that both sections are small, or both are the full residual row. Then $S\cap T$ is a relative support.
\end{lemma}
\begin{proof}
It is enough to generate the stabilizer of the intersection by the two stabilizers. On rows fixed entirely there is nothing to do. In each remaining row suppress $E$ and $S\cap T$ by working on their complement and write $D_1,D_2$ for the two disjoint small residual sections. Given a permissible permutation $g$, choose $u$ fixing $D_1$ so that $ug(D_1)$ misses $D_2$. This can be done by swapping the points of $g(D_1)\cap D_2$ with fresh points in a small set of equal size outside $D_1\cup D_2\cup g(D_1)$. The latter union is small, so its complement contains an infinite small subset. Then $ug|_{D_1}$ extends to a permissible permutation $v$ fixing $D_2$. For the extension, both its domain and range are small and disjoint from $D_2$. Take their union with a fresh infinite small set and biject its remaining domain to its remaining range, leaving everything else fixed. Thus $u$ and $v^{-1}ug$ fix $D_1$ and $v$ fixes $D_2$. The same three factors in each row assemble into three elements of the full product, proving the claim.
\end{proof}

\begin{lemma}\label{row}
On a countably infinite ground set $B$ with a tall proper ideal containing the finite sets let $H$ consist of all permutations with small moved set. In any $H$-action, a point with a countable orbit has a finite point support.
\end{lemma}
\begin{proof}
For every infinite small $R\subseteq B$, the embedded full group $\Sym(R)$ acts on a countable orbit. The stabilizer of the point in this group has countable index. By Dixon--Neumann--Thomas \cite[Theorem~1]{dnt}, it lies between the pointwise and setwise stabilizers of a finite $F_R\subseteq R$. This is the least finite support for the restricted action. Indeed any other finite support omitting $a\in F_R$ allows a transposition of $a$ with a fresh point outside both finite sets, contradicting setwise preservation of $F_R$.

If $R\subseteq R'$ are infinite small sets, then $F_{R'}\cap R$ supports the restricted action on $R$. Conversely, for any $a\in F_{R'}\cap R$, swapping it with any fresh point of $R\setminus(F_R\cup F_{R'})$ must move the point. Hence $a\in F_R$. Thus $F_R=F_{R'}\cap R$. Taking unions of two small sets gives compatibility for arbitrary $R,R'$. Consequently $F=\bigcup_R F_R$ satisfies $F\cap R=F_R$. If $F$ were infinite, tallness would give an infinite small subset $R\subseteq F$, an impossibility. Finally each element of $H$ moves points only in some infinite small $R$. If it fixes $F$, it fixes $F_R$ and hence the given point.
\end{proof}

\section{Which vertices can have complicated supports?}
\begin{lemma}\label{finite}
Let $V$ be an internal set supported by $E$ with no internally uncountable well-orderable subset. Each $x\in V$ has a least finite relative support.
\end{lemma}
\begin{proof}
Take a relative support $S$ for $x$, filling its finitely many exceptional sections outside $E$ up to whole residual rows. Let $D$ denote these finitely many new rows. In every other nonempty residual row $n$, take $2^n$ pairwise disjoint copies of $S_n$ inside one small reserve disjoint from $E_n\cup S_n$. This is possible even for infinite $S_n$, by tallness and finite partition of a denumerable reserve. For each binary real $r$ choose a permutation $g_r\in K$ moving $S_n$ to the copy indexed by $r|_n$, using swaps between the two copies. Leave the exceptional rows fixed. Every $g_rx$ lies in $V$ and is fixed by the stabilizer of one common allowed support consisting of $E$, the exceptional rows $D$ and all the reserves. Hence the set of all these images and its ground well-order belong to the model. By hypothesis the set is countable internally, hence externally. Two distinct reals $r,s$ give the same image.

The supports $g_rS$ and $g_sS$ of that image have disjoint sections at all sufficiently large $n$ outside $D$. Lemma~\ref{intersect} gives a relative support in finitely many rows. Moving back by $g_r^{-1}$ gives such a support for $x$. Enlarge it to those finitely many residual rows.

For any one of the latter rows, the orbit of $x$ under its residual row group is fixed pointwise by a common support (take $E$ and all the finitely many rows just obtained). Its ground well-order is internal and the orbit is countable. By Lemma~\ref{row}, the action in that row on $x$ has a finite support. Take the union of the resulting finitely many finite supports. It supports $x$ relative to $E$: factor any permutation fixing it into the finitely many individual row permutations and the factor fixing all these rows. Each fixes $x$. Finally two finite relative supports intersect to give a support, by Lemma~\ref{intersect}. A support of minimum finite cardinality is therefore included in every finite relative support. This proves leastness.
\end{proof}

\section{The partition theorem}
Let a graph on an uncountable internal $V$ have support $E$. If $V$ contains an internally uncountable well-orderable subset, restrict to it. The well-ordered version of $\BG$ holds in a permutation model \cite[Lemma~6.2]{ht}, e.g. by applying the ordinary partition theorem in the ground universe to a vertex subset fixed pointwise by one support. We may therefore assume the hypothesis of Lemma~\ref{finite}. Write $S_x$ for the least finite relative support of $x$. These supports satisfy $S_{gx}=gS_x$ for $g\in K$, by leastness.

Infinitely many rows occur in the union of all $S_x$. Otherwise, adding those finitely many whole rows to $E$ would fix every vertex, well-ordering $V$ internally. Choose externally vertices $x_i$, distinct rows $b_i$, and one point of $S_{x_i}$ in row $b_i$, for $i\in\omega$. We may successively take $b_i$ outside all rows of earlier chosen supports. In every nonempty residual row fix an infinite small reserve $R_n$ disjoint from $E_n$. Put $R=\bigcup_n R_n$, which is an allowed support. By independent permutations in $K$, place all the support points of each $x_i$ in $R$ except its one distinguished point. Place that point at $a_i\in A_{b_i}\setminus(E\cup R)$. This is possible with finitely many transpositions per chosen vertex. Write $y_i$ for the resulting vertex and $H=\Fix_G(E\cup R)$.

The orbit $V_i=Hy_i$ is an equivariant copy of $A_{b_i}\setminus(E\cup R)$. In fact an element of $H$ fixing $a_i$ fixes the support of $y_i$ pointwise. Conversely an element fixing $y_i$ preserves its least support setwise, and $a_i$ is its only point outside $E\cup R$. All transpositions of the residual row outside $R$ are available in $H$. These copies are disjoint, since equal vertices have equal least supports and the unique point outside $R$ would have to be in both distinguished rows. All pairs of distinct vertices within a $V_i$ have the same edge color, by the action on two-element subsets. All pairs between $V_i$ and $V_j$ have the same edge color for $i\ne j$, by independent action in their two rows.

An independent $V_i$ supplies a countably infinite independent set: the copy is infinite and the support $E\cup R$ enlarged by its one entire row fixes it pointwise. Suppose instead every $V_i$ is a clique. Apply the classical infinite Ramsey theorem in the ground universe to the cross-colors on $\omega$. If an infinite set of indices has independent cross-color, the corresponding $y_i$ form an internal countably infinite independent set. Their enumeration has support $E\cup R\cup\{a_i:i\in\omega\}$, which adds at most one point per row. If an infinite set $J$ of indices has clique cross-color, the union $\bigcup_{i\in J}V_i$ is an internal clique supported by $E\cup R$. It is internally uncountable. A support for a purported injection into $\omega$, enlarged by $E\cup R$, has only finitely many full rows. In some $b_i$, $i\in J$, its section is small; a transposition of two outside points fixes the support and interchanges their two vertices in $V_i$. It cannot preserve an injection of those vertices to natural numbers. This proves $\BG$ in the model.

The model also chooses from every well-orderable family of nonempty internally countable sets. More generally, use any well-orderable family $(V_i)$ in which no $V_i$ contains an uncountable well-orderable subset. A common support for the family and index order fixes each index, hence preserves each $V_i$. Choose one member of each $V_i$ externally and apply Lemma~\ref{finite} with that common parameter $E$. Move each of the finite supports into the same row reserves $R$. The choices thereby obtained all have support $E\cup R$ and their graph is internal.

\section{Transfer}
For a set $X$ write $|X|^-$ for the supremum of the well-ordered cardinals that inject into it. Let $\Psi(X)$ assert that every graph on $X$ either has countable vertex set, has a countably infinite independent set, or has an uncountable clique. This predicate is boundable. Its quantifiers can be restricted to $P^{\omega+20}(X)$ in the increasing power-set iteration: graph codes, subsets of vertices, $\omega$, and maps between these subsets and $\omega$ are all present there. Countability and its negation can be expressed by existence and nonexistence of the injection codes in this same stage. Thus
\[
 \forall X\bigl(|X|^-\le\omega\ \longrightarrow\ \Psi(X)\bigr)
\]
is an injectively boundable statement \cite[Section~5]{mad}.

Regularity of $\omega_1$ is boundable as a statement with existential ordinal witness $\kappa$. Require that $\kappa$ is an uncountable ordinal, every smaller ordinal is countable, and no map from a smaller ordinal to $\kappa$ is cofinal. All maps tested are bounded in the increasing hierarchy over $\kappa$ by $\omega+20$ steps. This assertion holds in the permutation model since the pure sets of the ground universe are preserved. Failure of $\mathrm{CUT}$ is also boundable: existentially specify a domain $A$ with a sequence of disjoint denumerable subsets covering $A$ and with no injection from $A$ into $\omega$. The enumeration or injection tests for the individual sets and the whole domain use maps within $P^{\omega+20}(A)$.

Pincus's theorem \cite[Theorem~3A3]{pincus} transfers a finite conjunction of injectively boundable statements, including ordinary boundable statements, from a permutation model to pure ZF. Apply it to the displayed restricted graph statement, regularity, and failure of $\mathrm{CUT}$. If a graph in the resulting ZF model has $|X|^->\omega$, it has a well-orderable subset of size $\aleph_1$. The well-orderable version of the partition theorem follows from regularity of $\aleph_1$ in ZF \cite[Proposition~3.3(1)]{bg}. Apply it to that subset; an infinite well-orderable independent set contains a countably infinite subset. For the other graphs the displayed statement already gives the conclusion. This proves Theorem~\ref{main}.

\paragraph{Relation to previous work.} Howard and Tachtsis \cite[Theorem~6.29 and Section~6.9]{ht} separate the graph principles from larger-cardinal union principles via a different model, in which the union of countably many denumerable sets remains countable. The present construction addresses precisely their remaining countable case and uses their distinction between the graph formulations. The symmetric-group small-index theorem and the permutation-to-ZF transfer are prior results. Independent review of the row construction, the transfer argument, and priority is pending.
\begin{thebibliography}{9}
\bibitem{ht} P. Howard and E. Tachtsis. \href{https://doi.org/10.4064/fm250514-20-7}{Three forms of the Erdős--Dushnik--Miller theorem}. Fund. Math. 272 (2026), 171--203. \href{https://arxiv.org/abs/2410.18229}{Author manuscript}, Section~8, Question~3.
\bibitem{bg} A. Banerjee and A. Gopaulsingh. \href{https://doi.org/10.4064/ba221221-6-6}{On the Erdős--Dushnik--Miller theorem without AC}. Bull. Polish Acad. Sci. Math. 71 (2023), 1--21.
\bibitem{dnt} J. D. Dixon, P. M. Neumann and S. Thomas. \href{https://doi.org/10.1112/blms/18.6.580}{Subgroups of small index in infinite symmetric groups}. Bull. London Math. Soc. 18 (1986), 580--586.
\bibitem{pincus} D. Pincus. \emph{Zermelo--Fraenkel consistency results by Fraenkel--Mostowski methods}. J. Symbolic Logic 37 (1972), 721--743.
\bibitem{mad} E. Tachtsis. \href{https://doi.org/10.4064/ba201014-22-1}{The Boolean prime ideal theorem does not imply the extension of almost disjoint families to MAD families}. Bull. Polish Acad. Sci. Math. 68 (2020), Section~5.
\end{thebibliography}
\end{document}
